\documentclass[12pt]{article}

\input{../../_assets/preambulo.tex}

\usepackage{algorithm}
\floatname{algorithm}{Algoritmo}
\usepackage{algpseudocode}

\algrenewcommand\algorithmicif{\textbf{Si}}
\algrenewcommand\algorithmicthen{\textbf{entonces}}
\algrenewcommand\algorithmicelse{\textbf{Si no}}
\algrenewcommand\algorithmicend{\textbf{Fin}}
\algrenewtext{For}[1]{\textbf{Para} #1 \textbf{hacer}}
\algrenewtext{EndFor}{\textbf{Fin Para}}
\algrenewtext{While}[1]{\textbf{Mientras} #1 \textbf{hacer}}
\algrenewtext{EndWhile}{\textbf{Fin Mientras}}
\algrenewtext{ElsIf}[1]{\textbf{Si no, si} #1 \textbf{entonces}}

\begin{document}

    % 1. Foto de fondo
    % 2. Título
    % 3. Encabezado Izquierdo
    % 4. Color de fondo
    % 5. Coord x del titulo
    % 6. Coord y del titulo
    % 7. Fecha

    
    \input{../../_assets/portada}
    \portadaExamen{ffccA4.jpg}{Modelos \\ Avanzados \\ de Computación \\Examen I}{Modelos Avanzados de Computación. Examen I}{MidnightBlue}{-8}{28}{2026}{José Manuel Sánchez Varbas}

    \begin{description}
        \item[Asignatura] Modelos Avanzados de Computación.
        \item[Curso Académico] 2024-25.
        \item[Grado] Grado en Ingeniería Informática.
        \item[Descripción] Ordinaria.
        \item[Fecha] 11 de Junio de 2025.  
    \end{description}
    \newpage


    % ------------------------------------

    \begin{ejercicio}[2.4 puntos]
        Determinar cuáles de los siguientes problemas son decidibles, semidecidibles o no semidecidibles, justificando las respuestas (se supone que las MTs tienen a $\{0,1\}$ como alfabeto de entrada):

        \begin{enumerate}
            \item Dada una MT $M$ de dos cintas y una palabra $u$ de longitud menor o igual a $100$, ¿es cierto que nunca usa más de $100$ casillas en la primera cinta cuando lee la palabra $u$?

            \item Dada una MT $M$, ¿es cierto que acepta todas las palabras de longitud menor o igual a $10$?

            \item Dadas dos MT, $M_1$ y $M_2$, ¿es cierto que toda palabra de longitud menor o igual a $100$ que es aceptada por $M_1$ lo es también por $M_2$ con un número de pasos menor o igual al empleado por $M_1$?

            \item Dadas dos MT, $M_1$ y $M_2$, y una palabra $u$, donde $M_2$ está garantizado que termina para todas las entradas, ¿es cierto que $M_1$ da más pasos que $M_2$ cuando ambas leen la palabra $u$?
        \end{enumerate}
    \end{ejercicio}

    \begin{ejercicio}[1.3 puntos]
        Describe de forma algorítmica una MT que acepte las palabras que son de la forma $uu$, donde $u \in \{0,1\}^{*}$ (es decir que sea una palabra repetida dos veces) en espacio logarítmico.
    \end{ejercicio}

    \begin{ejercicio}[2.4 puntos]
        Razona si son verdaderas o falsas las siguientes afirmaciones:

        \begin{enumerate}
            \item Una MT con una cinta ilimitada en ambas direcciones se puede simular siempre con una MT con una cinta ilimitada solo por la derecha y usando el mismo número de pasos.

            \item Para que una MT tenga una complejidad en espacio de $O(n^2)$ tiene que tener más de una cinta para no escribir en la cinta de entrada.

            \item Nuestro conocimiento actual es compatible con el hecho de que $\textbf{P} \neq \textbf{NP}$ y $\textbf{NP} = \textbf{CoNP}$.

            \item Está demostrado que la tesis de Church-Turing es cierta.
        \end{enumerate}
    \end{ejercicio}

    \begin{ejercicio}[1.3 puntos]
        Enuncia los problemas de la suma. ¿Qué sabes sobre su complejidad? Razona de forma breve la respuesta.
    \end{ejercicio}

    \begin{ejercicio}[1.3 puntos]
        Define la clase $\textbf{EXP}$, ¿qué relación existe con las clases $\textbf{PESPACIO}$ y $\textbf{P}$?
    \end{ejercicio}

    \begin{ejercicio}[1.3 puntos]
        Elige entre una de las siguientes opciones:

        \begin{itemize}
            \item Describe una reducción Turing de $\text{FSAT}(C)$ a $\text{SAT}(C)$.

            \item Sean los problemas de la existencia de un circuito hamiltoniano, $\text{HAMILTONIANO}(G)$, y el problema de buscar un circuito hamiltoniano 
            cuando este exista, $\text{FHAMILTONIANO}(G)$. Describe una reducción Turing de $\text{FHAMILTONIANO}(G)$ a $\text{HAMILTONIANO}(G)$.
        \end{itemize}
    \end{ejercicio}

    \newpage
    \setcounter{ejercicio}{0} % Reiniciar contador de ejercicios
    \noindent
    \textbf{Solución.}

    \begin{ejercicio}[2.4 puntos]
        Determinar cuáles de los siguientes problemas son decidibles, semidecidibles o no semidecidibles, justificando las respuestas (se supone que las MTs tienen a $\{0,1\}$ como alfabeto de entrada):

        \begin{enumerate}
            \item Dada una MT $M$ de dos cintas y una palabra $u$ de longitud menor o igual a $100$, ¿es cierto que nunca usa más de $100$ casillas 
            en la primera cinta cuando lee la palabra $u$? \\

            Definimos
            $$
            \hspace{-0.3cm} P_1=\{\langle M,u\rangle : |u|\leq 100 \text{ y nº casillas distintas 1º cinta visitadas por } M \text{ con entrada } u \leq 100\}  
            $$

            Veamos que no es semidecidible. Reducimos desde
            $$
            C\text{-PARADA}=\{\langle M,w\rangle : M \text{ no para con entrada } w\},
            $$
            que sabemos que no es semidecidible. Dada una instancia $\langle M,w\rangle$,
            construimos una MT de dos cintas $N$ y fijamos $u_0=\varepsilon$. \\

            La máquina $N$, con entrada $u_0$, simula $M$ con entrada $w$ usando solo la
            segunda cinta. Si la simulación de $M(w)$ para, entonces $N$ mueve el cabezal de
            la primera cinta hasta visitar 101 casillas distintas. \\

            Si $M$ no para con entrada $w$, entonces $N$ nunca visita más de 100 casillas
            distintas en la primera cinta, luego
            $$
            \langle N,u_0\rangle\in P_1.
            $$

            Si $M$ para con entrada $w$, entonces $N$ acaba visitando 101 casillas distintas
            de la primera cinta, luego
            $$
            \langle N,u_0\rangle\notin P_1.
            $$

            Por tanto,
            $$
            \langle M,w\rangle\in C\text{-PARADA}
            \Longleftrightarrow
            \langle N,u_0\rangle\in P_1.
            $$

            Luego $C\text{-PARADA}\propto P_1$. Como $C\text{-PARADA}$ no es
            semidecidible, $P_1$ tampoco es semidecidible. En particular, no es decidible.

            \newpage

            \item Dada una MT $M$, ¿es cierto que acepta todas las palabras de longitud menor o igual a $10$? \\
            
            Definimos
            $$
            P_2=\{\langle M\rangle : \forall w \in \{0,1\}^*,\ |w|\leq 10 \Rightarrow w\in L(M)\}.
            $$

            Este problema no es decidible por el Teorema de Rice, ya que la propiedad
            ``aceptar todas las palabras de longitud menor o igual que $10$'' depende solo del lenguaje $L(M)$,
            no de la máquina concreta. Además, es no trivial: la cumple una MT que acepta $\{0,1\}^*$, pero no la cumple
            una MT que no acepta ninguna palabra. \\

            Veamos que sí es semidecidible. Sea
            $$
            A=\{w\in\{0,1\}^* : |w|\leq 10\}.
            $$
            Como $A$ es finito, podemos escribir
            $$
            A=\{w_1,\dots,w_k\}.
            $$
            para cierto $k \in \mathbb{N}$. Construimos una MTND $N$ que funciona como sigue:

            \begin{algorithm}
            \caption{MTND $N$}
            \label{algorithm:todas-longitud-diez-nd}
                \begin{algorithmic}
                    \State \textbf{Entrada:} Una codificación $\langle M\rangle$ de una MT $M$.
                    \State Sea $A=\{w\in\{0,1\}^* : |w|\leq 10\}=\{w_1,\dots,w_k\}$.
                    \State Selecciona de forma no determinista un número natural $t$.
                    \For{$i=1,\dots,k$}
                        \State Simulamos $M$ con entrada $w_i$ durante $t$ pasos.
                        \If{$M$ no acepta $w_i$ en esos $t$ pasos}
                            \State $N$ rechaza.
                        \EndIf
                    \EndFor
                    \State $N$ acepta.
                \end{algorithmic}
            \end{algorithm}

            Si $M$ acepta todas las palabras de longitud menor o igual que $10$, entonces para cada $w_i$
            existe un número de pasos $t_i$ en el que $M$ acepta $w_i$. Por tanto, existe una rama de cómputo
            de $N$ que elige precisamente esos tiempos $t_1,\dots,t_k$, y esa rama acepta. \\

            Si $M$ no acepta alguna palabra $w_i$ de longitud menor o igual que $10$, entonces, elijamos los tiempos que elijamos,
            $M$ no aceptará esa palabra en esos pasos, y ninguna rama podrá aceptar. \\

            Por tanto, $P_2$ es semidecidible pero no decidible.

            \item Dadas dos MT, $M_1$ y $M_2$, ¿es cierto que toda palabra de longitud menor o igual a $100$ que es aceptada por $M_1$ lo es 
            también por $M_2$ con un número de pasos menor o igual al empleado por $M_1$? \\

            Definimos
            $$
            \hspace{-1.7cm}P_3=\{\langle M_1,M_2\rangle : \forall z,\ |z|\leq 100,\ 
            \text{si } M_1 \text{ acepta } z \text{ en } t_1 \text{ pasos, entonces } M_2
            \text{ acepta } z \text{ en } t_2\leq t_1 \text{ pasos}\}.
            $$

            Veamos que $P_3$ no es semidecidible. Reducimos desde
            $$
            C\text{-UNIVERSAL}(M,w)=\{\langle M,w\rangle : M \text{ no acepta } w\},
            $$
            que sabemos que no es semidecidible. Dada una instancia $\langle M,w\rangle$,
            construimos dos MT $N_1$ y $N_2$:

            \begin{algorithm}
            \caption{MT $N_1$}
            \label{algorithm:N1-p3-corto}
            \begin{algorithmic}
            \State \textbf{Entrada:} Una palabra $x\in\{0,1\}^*$.
            \If{$x\neq \varepsilon$}
                \State Rechazar.
            \EndIf
            \State Simular $M$ con entrada $w$.
            \If{$M$ acepta $w$}
                \State Aceptar.
            \EndIf
            \end{algorithmic}
            \end{algorithm}

            \begin{algorithm}
            \caption{MT $N_2$}
            \label{algorithm:N2-p3-corto}
            \begin{algorithmic}
            \State \textbf{Entrada:} Una palabra $x\in\{0,1\}^*$.
            \State Rechazar.
            \end{algorithmic}
            \end{algorithm}

            Si $M$ no acepta $w$, entonces $N_1$ no acepta ninguna palabra. Por tanto,
            la condición de $P_3$ se cumple trivialmente, ya que no hay palabras aceptadas
            por $N_1$. \\

            Si $M$ acepta $w$, entonces $N_1$ acepta $\varepsilon$ en cierto número $t_1$
            de pasos, pero $N_2$ no acepta $\varepsilon$ en ningún número de pasos. Por
            tanto, la condición de $P_3$ falla. \\

            Así,
            $$
            \langle M,w\rangle\in C\text{-UNIVERSAL}
            \Longleftrightarrow
            \langle N_1,N_2\rangle\in P_3.
            $$

            Luego $C\text{-UNIVERSAL}\propto P_3$. Como $C\text{-UNIVERSAL}$ no es
            semidecidible, $P_3$ no es semidecidible. En particular, $P_3$ no es decidible.

            \newpage

            \item Dadas dos MT, $M_1$ y $M_2$, y una palabra $u$, donde $M_2$ está garantizado que termina para todas las entradas, 
            ¿es cierto que $M_1$ da más pasos que $M_2$ cuando ambas leen la palabra $u$? \\

            Definimos
            $$
            P_4=\{\langle M_1,M_2,u\rangle : M_1 \text{ da más pasos que } M_2
            \text{ con entrada } u\}.
            $$

            Este problema es decidible. Como $M_2$ está garantizado que termina para todas
            las entradas, en particular termina con entrada $u$. Por tanto, podemos simular
            $M_2(u)$ hasta que pare y calcular su número de pasos, digamos $t_2$. \\

            Después simulamos $M_1(u)$ durante exactamente $t_2$ pasos.

            \begin{algorithm}
            \caption{$M_{P_4}$}
            \label{algorithm:p4-decidible}
            \begin{algorithmic}
            \State \textbf{Entrada:} Una codificación $\langle M_1,M_2,u\rangle$.
            \State Simular $M_2$ con entrada $u$ hasta que termine.
            \State Sea $t_2$ el número de pasos dados por $M_2$.
            \State Simular $M_1$ con entrada $u$ durante exactamente $t_2$ pasos.
            \If{$M_1$ no ha terminado en ningún momento de esos $t_2$ pasos}
                \State Aceptar.
            \Else
                \State Rechazar.
            \EndIf
            \end{algorithmic}
            \end{algorithm}

            El algoritmo siempre termina, porque $M_2(u)$ termina por hipótesis y luego
            solo se simula $M_1(u)$ durante un número finito de pasos. Además, si $M_1$
            no ha terminado en los primeros $t_2$ pasos, entonces da más pasos que $M_2$;
            mientras que si ha terminado en algún instante $t_1\leq t_2$, no da más pasos
            que $M_2$. \\

            Por tanto, $P_4$ es decidible.
        \end{enumerate}
    \end{ejercicio}

    \newpage

    \begin{ejercicio}[1.3 puntos]
        Describe de forma algorítmica una MT que acepte las palabras que son de la forma $uu$, donde $u \in \{0,1\}^{*}$ 
        (es decir que sea una palabra repetida dos veces) en espacio logarítmico. \\

        Sea $x$ la palabra de entrada y sea $n=|x|$. Queremos decidir si existe
        $u\in\{0,1\}^*$ tal que $x=uu$. Construimos una MT determinista con cinta
        de entrada de solo lectura y una cinta de trabajo.

        \begin{algorithm}
        \caption{MT para $L=\{uu : u\in\{0,1\}^*\}$}
        \label{algorithm:uu-log}
        \begin{algorithmic}
        \State \textbf{Entrada:} Una palabra $x\in\{0,1\}^*$.
        \State Calcular $n=|x|$ en binario usando un contador.
        \If{$n$ es impar}
            \State Rechazar.
        \EndIf
        \State Calcular $m=n/2$ eliminando el último bit $0$ de la representación binaria de $n$.
        \For{$i=1,\dots,m$}
            \State Mover el cabezal de entrada hasta la posición $i$ y guardar el símbolo leído en una variable $a$.
            \State Mover el cabezal de entrada hasta la posición $i+m$ y guardar el símbolo leído en una variable $b$.
            \If{$a\neq b$}
                \State Rechazar.
            \EndIf
        \EndFor
        \State Aceptar.
        \end{algorithmic}
        \end{algorithm}

        La máquina acepta exactamente las palabras de la forma $uu$, ya que compara,
        posición a posición, la primera mitad de $x$ con la segunda mitad. Si todas las
        posiciones coinciden, entonces $x=uu$, con $u$ igual a la primera mitad de $x$;
        si alguna no coincide, entonces $x$ no puede ser de esa forma. \\

        Veamos el espacio usado. La máquina solo necesita almacenar los contadores
        $n$, $m$ e $i$, además de dos símbolos $a,b\in\{0,1\}$. Como estos contadores
        toman valores menores o iguales que $n$, cada uno ocupa $O(\log n)$ casillas.
        Además, calcular $m=n/2$ no requiere realizar una división general: como ya se
        ha comprobado que $n$ es par, su representación binaria acaba en $0$, y dividir
        por $2$ consiste simplemente en eliminar ese último bit. Por tanto, ese cálculo
        no aumenta el orden de espacio usado. Así, el espacio total es $O(\log n)$. \\

        Luego
        $$
        L=\{uu : u\in\{0,1\}^*\}
        $$
        se reconoce en espacio logarítmico.
    \end{ejercicio}

    \newpage

    \begin{ejercicio}[2.4 puntos]
        Razona si son verdaderas o falsas las siguientes afirmaciones:
        \begin{enumerate}
            \item Una MT con una cinta ilimitada en ambas direcciones se puede simular siempre con una MT con una cinta ilimitada solo por la derecha y 
            usando el mismo número de pasos. \\

            \boxed{\text{Falsa}}. Una MT con cinta ilimitada en ambas direcciones sí puede simularse con una cinta limitada solo por la derecha usando dos pistas, 
            colocando en una pista las posiciones $0,1,2,\dots$ y en otra las posiciones $-1,-2,\dots$. Sin embargo, no se garantiza que use exactamente el mismo 
            número de pasos: la construcción introduce estados auxiliares para preparar la cinta y gestionar qué pista está activa. Por tanto, la simulación existe, 
            pero no necesariamente conserva exactamente el número de pasos.

            \item Para que una MT tenga una complejidad en espacio de $O(n^2)$ tiene que tener más de una cinta para no escribir en la cinta de entrada. \\
            
            \boxed{\text{Falsa}}. Como contraejemplo, consideramos una MT de una sola cinta que, con entrada
            $x\in\{0,1\}^*$ de longitud $n$, escribe $n$ copias de $x$ a la derecha de la entrada.

            De forma algorítmica:

            \begin{algorithm}
            \caption{MT de una cinta que usa espacio $O(n^2)$}
            \begin{algorithmic}
            \State \textbf{Entrada:} Una palabra $x\in\{0,1\}^*$.
            \State Recorre la entrada de izquierda a derecha.
            \For{cada símbolo de la entrada}
                \State Vuelve al principio de la entrada.
                \State Copia toda la palabra $x$ al final de la zona ya escrita.
            \EndFor
            \State Acepta.
            \end{algorithmic}
            \end{algorithm}

            Como el bucle exterior se ejecuta una vez por cada símbolo de $x$, se ejecuta $n$ veces.
            En cada iteración copia una palabra de longitud $n$. Por tanto, al final ha escrito $n$
            copias de $x$, es decir, ha usado al menos $n^2$ casillas adicionales. Esta MT tiene una sola cinta,
            luego no es necesario tener más de una cinta para tener complejidad espacial $O(n^2)$.

            \item Nuestro conocimiento actual es compatible con el hecho de que $\textbf{P} \neq \textbf{NP}$ y $\textbf{NP} = \textbf{CoNP}$. \\
            
            \boxed{\text{Verdadera}}. Es compatible con el conocimiento actual que $\textbf{P}\neq\textbf{NP}$ y, a la vez, $\textbf{NP}=\textbf{CoNP}$. 
            Sabemos que $\textbf{P}\subseteq\textbf{NP}$ y $\textbf{P}\subseteq\textbf{CoNP}$, pero no está demostrado ni que $\textbf{P}=\textbf{NP}$, ni 
            que $\textbf{NP}=\textbf{CoNP}$, ni que $\textbf{NP}\neq\textbf{CoNP}$.

            \newpage

            \item Está demostrado que la tesis de Church-Turing es cierta. \\
            
            \boxed{\text{Falsa}}. ``Es una tesis que no puede ser comprobada matemáticamente si no se da una definición formal de efectivamente calculable, 
            pero eso es precisamente lo que intenta hacer la tesis de Church Turing. Es una tesis que es aceptada como cierta. Aunque consideremos modelos
            de computación cuántica, no permitirían hacer cálculos más allá de los realizados por una MT (aunque quizá en menos tiempo que una MT).''
        \end{enumerate}
    \end{ejercicio}

    \begin{ejercicio}[1.3 puntos]
        Enuncia los problemas de la suma. ¿Qué sabes sobre su complejidad? Razona de forma breve la respuesta. \\

        El problema de la SUMA se enuncia como sigue: dados un conjunto finito $A$, una función de tamaños
        $
        s:A\to \mathbb{N},
        $
        y un entero $B$, determinar si existe un subconjunto $A'\subseteq A$ tal que
        $$
        \sum_{a\in A'} s(a)=B.
        $$

        Este problema está en $\textbf{NP}$, porque basta elegir no deterministamente el subconjunto
        $A'$ y comprobar en tiempo polinómico si la suma de sus tamaños es exactamente $B$. \\

        Además, SUMA es $\textbf{NP}$-completo. Según lo que se ha visto en teoría se ha demostrado reduciendo
        ACTRI a SUMA: para cada tripleta se construye un número, y se fija un valor $B$ de forma
        que elegir números que suman $B$ equivale exactamente a elegir un acoplamiento de tripletas.
        Como ACTRI es $\textbf{NP}$-completo y ACTRI $\propto$ SUMA, se concluye que SUMA es
        $\textbf{NP}$-completo. \\

        También aparece relacionado el problema de la \textbf{PARTICIÓN}: dado un conjunto $C$ con
        tamaños $s:C\to\mathbb{N}$, determinar si existe $C'\subseteq C$ tal que
        $$
        \sum_{a\in C'}s(a)=\sum_{a\in C\setminus C'}s(a).
        $$
        En teoría hemos visto cómo reducir SUMA a PARTICIÓN, por lo que PARTICIÓN también es
        $\textbf{NP}$-completo.
    \end{ejercicio}

    \newpage

    \begin{ejercicio}[1.3 puntos]
        Define la clase $\textbf{EXP}$, ¿qué relación existe con las clases $\textbf{PESPACIO}$ y $\textbf{P}$? \\

        La clase $\textbf{EXP}$ es la clase de los lenguajes decidibles en tiempo
        exponencial determinista:
        $$
        \textbf{EXP}=\bigcup_{j>0}\textbf{TIEMPO}(2^{n^j}).
        $$

        Según se ha visto en teoría:
        $$
        \textbf{P}=\bigcup_{j>0}\textbf{TIEMPO}(n^j),
        \qquad
        \textbf{PESPACIO}=\bigcup_{j>0}\textbf{ESPACIO}(n^j).
        $$

        La relación entre ellas es
        $$
        \textbf{P}\subseteq \textbf{PESPACIO}\subseteq \textbf{EXP}.
        $$

        En efecto, si un algoritmo trabaja en tiempo polinómico, entonces no puede usar
        más que espacio polinómico, luego
        $$
        \textbf{P}\subseteq \textbf{PESPACIO}.
        $$
        Por otro lado, si una MT usa espacio polinómico, entonces solo tiene
        exponencialmente muchas configuraciones posibles; por tanto, se puede decidir en
        tiempo exponencial si acepta o no, luego
        $$
        \textbf{PESPACIO}\subseteq \textbf{EXP}.
        $$

        Además, por el teorema de jerarquía en tiempo se sabe que
        $$
        \textbf{P}\neq \textbf{EXP}.
        $$
        Por tanto, al menos una de las dos inclusiones anteriores es estricta:
        $$
        \textbf{P}\subsetneq \textbf{PESPACIO}
        \qquad\text{o}\qquad
        \textbf{PESPACIO}\subsetneq \textbf{EXP},
        $$
        pero no se sabe cuál.
    \end{ejercicio}

    \newpage

    \begin{ejercicio}[1.3 puntos]
        Elige entre una de las siguientes opciones:

        \begin{itemize}
            \item Describe una reducción Turing de $\text{FSAT}(C)$ a $\text{SAT}(C)$. \\
            
            El problema $\textbf{SAT}(C)$ decide si un conjunto de cláusulas $C$ es satisfacible. El problema
            $\textbf{FSAT}(C)$ consiste en encontrar una asignación satisfactoria para $C$, si existe. \\

            Suponemos que tenemos un oráculo para $\textbf{SAT}$. Construimos una reducción Turing para
            $\textbf{FSAT}$:

            \begin{algorithm}
            \caption{Reducción Turing de $\textbf{FSAT}$ a $\textbf{SAT}$}
            \begin{algorithmic}
            \State \textbf{Entrada:} Un conjunto de cláusulas $C$ con variables $x_1,\dots,x_n$.
            \If{$\textbf{SAT}(C)$ responde NO}
                \State Responder que no existe asignación satisfactoria.
            \EndIf
            \For{$i=1,\dots,n$}
                \State Sea $C_0=C\cup\{\neg x_i\}$.
                \If{$\textbf{SAT}(C_0)$ responde SI}
                    \State Fijar $x_i=0$ y sustituirlo en $C$.
                \Else
                    \State Fijar $x_i=1$ y sustituirlo en $C$.
                \EndIf
            \EndFor
            \State Devolver la asignación obtenida.
            \end{algorithmic}
            \end{algorithm}

            El algoritmo hace a lo sumo $n+1$ llamadas al oráculo $\textbf{SAT}$. En cada paso mantiene una
            fórmula satisfacible: si al fijar $x_i=0$ sigue siendo satisfacible, tomamos $x_i=0$; si no,
            como la fórmula anterior era satisfacible, necesariamente existe una solución con $x_i=1$.
            Por tanto, al final se obtiene una asignación satisfactoria de $C$. Luego
            $$
            \textbf{FSAT} \propto_T \textbf{SAT}.
            $$

            \newpage

            \item Sean los problemas de la existencia de un circuito hamiltoniano, $\text{HAMILTONIANO}(G)$, y el problema de buscar un circuito hamiltoniano 
            cuando este exista, $\text{FHAMILTONIANO}(G)$. Describe una reducción Turing de $\text{FHAMILTONIANO}(G)$ a $\text{HAMILTONIANO}(G)$. \\

            Suponemos que tenemos un oráculo para $\textbf{HAMILTONIANO}$. Construimos una reducción
            Turing para $\textbf{FHAMILTONIANO}$: 

            \begin{algorithm}
            \caption{Reducción Turing de $\textbf{FHAMILTONIANO}$ a $\textbf{HAMILTONIANO}$}
            \begin{algorithmic}
            \State \textbf{Entrada:} Un grafo no dirigido $G=(V,E)$.
            \If{$\textbf{HAMILTONIANO}(G)$ responde NO}
                \State Responder que no existe circuito hamiltoniano.
            \EndIf
            \State $G'\gets G$.
            \For{cada arista $e\in E$}
                \State Sea $G_e=G'-e$.
                \If{$\textbf{HAMILTONIANO}(G_e)$ responde SI}
                    \State $G'\gets G_e$.
                \EndIf
            \EndFor
            \State Devolver el circuito formado por las aristas de $G'$.
            \end{algorithmic}
            \end{algorithm}

            El algoritmo elimina una arista siempre que, al quitarla, el grafo sigue teniendo algún circuito
            hamiltoniano. Por tanto, mantiene invariante que $G'$ tiene un circuito hamiltoniano. Al final,
            ninguna arista sobrante puede eliminarse sin destruir todos los circuitos hamiltonianos; en
            particular, las aristas que quedan contienen exactamente un circuito hamiltoniano. Ese circuito
            se devuelve. Luego
            $$
            \textbf{FHAMILTONIANO} \propto_T \textbf{HAMILTONIANO}.
            $$

        \end{itemize}
    \end{ejercicio}

\end{document}
